\IfFileExists{stacks-project.cls}{%
\documentclass{stacks-project}
}{%
\documentclass{amsart}
}

% For dealing with references we use the comment environment
\usepackage{verbatim}
\newenvironment{reference}{\comment}{\endcomment}
%\newenvironment{reference}{}{}
\newenvironment{slogan}{\comment}{\endcomment}
\newenvironment{history}{\comment}{\endcomment}

% For commutative diagrams we use Xy-pic
\usepackage[all]{xy}

% We use 2cell for 2-commutative diagrams.
\xyoption{2cell}
\UseAllTwocells

% We use multicol for the list of chapters between chapters
\usepackage{multicol}

% This is generally recommended for better output
\usepackage{lmodern}
\usepackage[T1]{fontenc}

% For cross-file-references
\usepackage{xr-hyper}

% Package for hypertext links:
\usepackage{hyperref}

% For any local file, say "hello.tex" you want to link to please
% use \externaldocument[hello-]{hello}
\externaldocument[introduction-]{introduction}
\externaldocument[conventions-]{conventions}
\externaldocument[sets-]{sets}
\externaldocument[categories-]{categories}
\externaldocument[topology-]{topology}
\externaldocument[sheaves-]{sheaves}
\externaldocument[sites-]{sites}
\externaldocument[stacks-]{stacks}
\externaldocument[fields-]{fields}
\externaldocument[algebra-]{algebra}
\externaldocument[brauer-]{brauer}
\externaldocument[homology-]{homology}
\externaldocument[derived-]{derived}
\externaldocument[simplicial-]{simplicial}
\externaldocument[more-algebra-]{more-algebra}
\externaldocument[smoothing-]{smoothing}
\externaldocument[modules-]{modules}
\externaldocument[sites-modules-]{sites-modules}
\externaldocument[injectives-]{injectives}
\externaldocument[cohomology-]{cohomology}
\externaldocument[sites-cohomology-]{sites-cohomology}
\externaldocument[dga-]{dga}
\externaldocument[dpa-]{dpa}
\externaldocument[sdga-]{sdga}
\externaldocument[hypercovering-]{hypercovering}
\externaldocument[schemes-]{schemes}
\externaldocument[constructions-]{constructions}
\externaldocument[properties-]{properties}
\externaldocument[morphisms-]{morphisms}
\externaldocument[coherent-]{coherent}
\externaldocument[divisors-]{divisors}
\externaldocument[limits-]{limits}
\externaldocument[varieties-]{varieties}
\externaldocument[topologies-]{topologies}
\externaldocument[descent-]{descent}
\externaldocument[perfect-]{perfect}
\externaldocument[more-morphisms-]{more-morphisms}
\externaldocument[flat-]{flat}
\externaldocument[groupoids-]{groupoids}
\externaldocument[more-groupoids-]{more-groupoids}
\externaldocument[etale-]{etale}
\externaldocument[chow-]{chow}
\externaldocument[intersection-]{intersection}
\externaldocument[pic-]{pic}
\externaldocument[weil-]{weil}
\externaldocument[adequate-]{adequate}
\externaldocument[dualizing-]{dualizing}
\externaldocument[duality-]{duality}
\externaldocument[discriminant-]{discriminant}
\externaldocument[derham-]{derham}
\externaldocument[local-cohomology-]{local-cohomology}
\externaldocument[algebraization-]{algebraization}
\externaldocument[curves-]{curves}
\externaldocument[resolve-]{resolve}
\externaldocument[models-]{models}
\externaldocument[functors-]{functors}
\externaldocument[equiv-]{equiv}
\externaldocument[pione-]{pione}
\externaldocument[etale-cohomology-]{etale-cohomology}
\externaldocument[proetale-]{proetale}
\externaldocument[relative-cycles-]{relative-cycles}
\externaldocument[more-etale-]{more-etale}
\externaldocument[trace-]{trace}
\externaldocument[crystalline-]{crystalline}
\externaldocument[spaces-]{spaces}
\externaldocument[spaces-properties-]{spaces-properties}
\externaldocument[spaces-morphisms-]{spaces-morphisms}
\externaldocument[decent-spaces-]{decent-spaces}
\externaldocument[spaces-cohomology-]{spaces-cohomology}
\externaldocument[spaces-limits-]{spaces-limits}
\externaldocument[spaces-divisors-]{spaces-divisors}
\externaldocument[spaces-over-fields-]{spaces-over-fields}
\externaldocument[spaces-topologies-]{spaces-topologies}
\externaldocument[spaces-descent-]{spaces-descent}
\externaldocument[spaces-perfect-]{spaces-perfect}
\externaldocument[spaces-more-morphisms-]{spaces-more-morphisms}
\externaldocument[spaces-flat-]{spaces-flat}
\externaldocument[spaces-groupoids-]{spaces-groupoids}
\externaldocument[spaces-more-groupoids-]{spaces-more-groupoids}
\externaldocument[bootstrap-]{bootstrap}
\externaldocument[spaces-pushouts-]{spaces-pushouts}
\externaldocument[spaces-chow-]{spaces-chow}
\externaldocument[groupoids-quotients-]{groupoids-quotients}
\externaldocument[spaces-more-cohomology-]{spaces-more-cohomology}
\externaldocument[spaces-simplicial-]{spaces-simplicial}
\externaldocument[spaces-duality-]{spaces-duality}
\externaldocument[formal-spaces-]{formal-spaces}
\externaldocument[restricted-]{restricted}
\externaldocument[spaces-resolve-]{spaces-resolve}
\externaldocument[formal-defos-]{formal-defos}
\externaldocument[defos-]{defos}
\externaldocument[cotangent-]{cotangent}
\externaldocument[examples-defos-]{examples-defos}
\externaldocument[algebraic-]{algebraic}
\externaldocument[examples-stacks-]{examples-stacks}
\externaldocument[stacks-sheaves-]{stacks-sheaves}
\externaldocument[criteria-]{criteria}
\externaldocument[artin-]{artin}
\externaldocument[quot-]{quot}
\externaldocument[stacks-properties-]{stacks-properties}
\externaldocument[stacks-morphisms-]{stacks-morphisms}
\externaldocument[stacks-limits-]{stacks-limits}
\externaldocument[stacks-cohomology-]{stacks-cohomology}
\externaldocument[stacks-perfect-]{stacks-perfect}
\externaldocument[stacks-introduction-]{stacks-introduction}
\externaldocument[stacks-more-morphisms-]{stacks-more-morphisms}
\externaldocument[stacks-geometry-]{stacks-geometry}
\externaldocument[moduli-]{moduli}
\externaldocument[moduli-curves-]{moduli-curves}
\externaldocument[examples-]{examples}
\externaldocument[exercises-]{exercises}
\externaldocument[guide-]{guide}
\externaldocument[desirables-]{desirables}
\externaldocument[coding-]{coding}
\externaldocument[obsolete-]{obsolete}
\externaldocument[fdl-]{fdl}
\externaldocument[index-]{index}

% Theorem environments.
%
\theoremstyle{plain}
\newtheorem{theorem}[subsection]{Theorem}
\newtheorem{proposition}[subsection]{Proposition}
\newtheorem{lemma}[subsection]{Lemma}

\theoremstyle{definition}
\newtheorem{definition}[subsection]{Definition}
\newtheorem{example}[subsection]{Example}
\newtheorem{exercise}[subsection]{Exercise}
\newtheorem{situation}[subsection]{Situation}

\theoremstyle{remark}
\newtheorem{remark}[subsection]{Remark}
\newtheorem{remarks}[subsection]{Remarks}

\numberwithin{equation}{subsection}

% Macros
%
\def\lim{\mathop{\mathrm{lim}}\nolimits}
\def\colim{\mathop{\mathrm{colim}}\nolimits}
\def\Spec{\mathop{\mathrm{Spec}}}
\def\Hom{\mathop{\mathrm{Hom}}\nolimits}
\def\Ext{\mathop{\mathrm{Ext}}\nolimits}
\def\SheafHom{\mathop{\mathcal{H}\!\mathit{om}}\nolimits}
\def\SheafExt{\mathop{\mathcal{E}\!\mathit{xt}}\nolimits}
\def\Sch{\mathit{Sch}}
\def\Mor{\mathop{\mathrm{Mor}}\nolimits}
\def\Ob{\mathop{\mathrm{Ob}}\nolimits}
\def\Sh{\mathop{\mathit{Sh}}\nolimits}
\def\NL{\mathop{N\!L}\nolimits}
\def\CH{\mathop{\mathrm{CH}}\nolimits}
\def\proetale{{pro\text{-}\acute{e}tale}}
\def\etale{{\acute{e}tale}}
\def\QCoh{\mathit{QCoh}}
\def\Ker{\mathop{\mathrm{Ker}}}
\def\Im{\mathop{\mathrm{Im}}}
\def\Coker{\mathop{\mathrm{Coker}}}
\def\Coim{\mathop{\mathrm{Coim}}}

% Boxtimes
%
\DeclareMathSymbol{\boxtimes}{\mathbin}{AMSa}{"02}

%
% Macros for moduli stacks/spaces
%
\def\QCohstack{\mathcal{QC}\!\mathit{oh}}
\def\Cohstack{\mathcal{C}\!\mathit{oh}}
\def\Spacesstack{\mathcal{S}\!\mathit{paces}}
\def\Quotfunctor{\mathrm{Quot}}
\def\Hilbfunctor{\mathrm{Hilb}}
\def\Curvesstack{\mathcal{C}\!\mathit{urves}}
\def\Polarizedstack{\mathcal{P}\!\mathit{olarized}}
\def\Complexesstack{\mathcal{C}\!\mathit{omplexes}}
% \Pic is the operator that assigns to X its picard group, usage \Pic(X)
% \Picardstack_{X/B} denotes the Picard stack of X over B
% \Picardfunctor_{X/B} denotes the Picard functor of X over B
\def\Pic{\mathop{\mathrm{Pic}}\nolimits}
\def\Picardstack{\mathcal{P}\!\mathit{ic}}
\def\Picardfunctor{\mathrm{Pic}}
\def\Deformationcategory{\mathcal{D}\!\mathit{ef}}
\usepackage{cleveref}
\crefname{originalreference}{original reference}{original references}
\Crefname{originalreference}{Original reference}{Original references}
\makeatletter
\newlabel{algebra-lemma-colimit-syntomic}{{07BV}{source}{}{sourcepointer-1}{}}
\newlabel{algebra-lemma-colimit-syntomic@cref}{{[originalreference][0][]07BV}{[1][0][]source}{}{}{}}
\makeatother
\crefname{originalreference}{original reference}{original references}
\Crefname{originalreference}{Original reference}{Original references}
\makeatletter
\newlabel{algebra-lemma-flat-over-regular}{{07DY}{source}{}{sourcepointer-2}{}}
\newlabel{algebra-lemma-flat-over-regular@cref}{{[originalreference][0][]07DY}{[1][0][]source}{}{}{}}
\makeatother
\crefname{originalreference}{original reference}{original references}
\Crefname{originalreference}{Original reference}{Original references}
\makeatletter
\newlabel{algebra-lemma-when-colimit}{{07C3}{source}{}{sourcepointer-3}{}}
\newlabel{algebra-lemma-when-colimit@cref}{{[originalreference][0][]07C3}{[1][0][]source}{}{}{}}
\makeatother
\crefname{originallemma}{original lemma}{original lemmas}
\Crefname{originallemma}{Original lemma}{Original lemmas}
\makeatletter
\newlabel{lemma-compare-standard}{{07EZ}{source}{}{sourcepointer-4}{}}
\newlabel{lemma-compare-standard@cref}{{[originallemma][0][]07EZ}{[1][0][]source}{}{}{}}
\makeatother
\crefname{originallemma}{original lemma}{original lemmas}
\Crefname{originallemma}{Original lemma}{Original lemmas}
\makeatletter
\newlabel{lemma-delocalize-height-zero}{{07FA}{source}{}{sourcepointer-5}{}}
\newlabel{lemma-delocalize-height-zero@cref}{{[originallemma][0][]07FA}{[1][0][]source}{}{}{}}
\makeatother
\crefname{originallemma}{original lemma}{original lemmas}
\Crefname{originallemma}{Original lemma}{Original lemmas}
\makeatletter
\newlabel{lemma-find-sequence}{{07FD}{source}{}{sourcepointer-6}{}}
\newlabel{lemma-find-sequence@cref}{{[originallemma][0][]07FD}{[1][0][]source}{}{}{}}
\makeatother
\crefname{originallemma}{original lemma}{original lemmas}
\Crefname{originallemma}{Original lemma}{Original lemmas}
\makeatletter
\newlabel{lemma-improve-presentation}{{07CE}{source}{}{sourcepointer-7}{}}
\newlabel{lemma-improve-presentation@cref}{{[originallemma][0][]07CE}{[1][0][]source}{}{}{}}
\makeatother
\crefname{originallemma}{original lemma}{original lemmas}
\Crefname{originallemma}{Original lemma}{Original lemmas}
\makeatletter
\newlabel{lemma-lift-solution}{{07F8}{source}{}{sourcepointer-8}{}}
\newlabel{lemma-lift-solution@cref}{{[originallemma][0][]07F8}{[1][0][]source}{}{}{}}
\makeatother
\crefname{originalproposition}{original proposition}{original propositions}
\Crefname{originalproposition}{Original proposition}{Original propositions}
\makeatletter
\newlabel{proposition-lift}{{07CM}{source}{}{sourcepointer-9}{}}
\newlabel{proposition-lift@cref}{{[originalproposition][0][]07CM}{[1][0][]source}{}{}{}}
\makeatother
\crefname{originalreference}{original reference}{original references}
\Crefname{originalreference}{Original reference}{Original references}
\makeatletter
\newlabel{situation-local}{{07F7}{source}{}{sourcepointer-10}{}}
\newlabel{situation-local@cref}{{[originalreference][0][]07F7}{[1][0][]source}{}{}{}}
\makeatother
\title{Resolving at a prime with regular local target and separable residue extension}
\author{The Stacks Project Authors}
\date{stacks-project:smoothing.tex:local-snapshot-20261012; GNU FDL 1.2 or later}
\begin{document}
\maketitle
\noindent\textit{Editorial note.} Inputs lemma-improve-presentation (07CE) and lemma-compare-standard (07EZ) are published entries used as standing inputs by their display numbers; not proven on the author's behalf.\par\medskip
\noindent\textit{Allowed standard prerequisites (editorial).}\begin{itemize}
\item Presentation improvement with retraction. lemma-improve-presentation (same chapter, PUBLISHED 07CE). Standardizes the auxiliary algebra B and freeness of the differential module.
\item Comparison of standardness conditions. lemma-compare-standard (same chapter, PUBLISHED 07EZ). Promotes smoothness with free differentials to strict standardness of a power x\_i\textasciicircum{}c.
\item Parameter sequence with controlled annihilators. lemma-find-sequence (carded L26815). Supplies pi\_1, ..., pi\_d with pi\_i\textasciicircum{}n in h\_A and the annihilator equalities for e = 8c.
\item Removing strictly standard elements from a resolution. lemma-lift-solution (carded L26812). Reduces the resolution to the quotient by the powers x\_i\textasciicircum{}e.
\item Height-zero delocalization. lemma-delocalize-height-zero (carded L26813). Reduces further to the localization at q.
\item Ind-smoothness under nilpotent thickenings. proposition-lift (carded L26807). Identifies the flat Artinian quotient map as a filtered colimit of smooth algebras.
\item Colimit of syntomic maps. Algebra, lemma colimit-syntomic (cross-chapter, named). Recognizes the flat map with separable residue extension as a filtered colimit of syntomic algebras.
\item Colimit criterion for ind-smooth maps. Algebra, lemma when-colimit (cross-chapter, named). Extracts the solution from the colimit presentation.
\end{itemize}\medskip
\ifdefined\hypertarget\hypertarget{sourcepointer-1}{}\fi
\noindent\textit{Source reference.} Original reference 07BV: focused-seeds/geometry/S9-smo-resolve-special/ref-tags.txt, source lines 5--5. Navigation pointer only: the counted proof cites this statement by \textbackslash{}ref as a standing input; the deduction the counted proof has to make is carried out inside the counted span, and no part of that deduction is delegated to this pointer.\par\medskip
\ifdefined\hypertarget\hypertarget{sourcepointer-2}{}\fi
\noindent\textit{Source reference.} Original reference 07DY: focused-seeds/geometry/S9-smo-resolve-special/ref-tags.txt, source lines 6--6. Navigation pointer only: the counted proof cites this statement by \textbackslash{}ref as a standing input; the deduction the counted proof has to make is carried out inside the counted span, and no part of that deduction is delegated to this pointer.\par\medskip
\ifdefined\hypertarget\hypertarget{sourcepointer-3}{}\fi
\noindent\textit{Source reference.} Original reference 07C3: focused-seeds/geometry/S9-smo-resolve-special/ref-tags.txt, source lines 7--7. Navigation pointer only: the counted proof cites this statement by \textbackslash{}ref as a standing input; the deduction the counted proof has to make is carried out inside the counted span, and no part of that deduction is delegated to this pointer.\par\medskip
\ifdefined\hypertarget\hypertarget{sourcepointer-4}{}\fi
\noindent\textit{Source reference.} Original lemma 07EZ: sources/stacks-smoothing/pinned-original-source.tex, source lines 936--936. Navigation pointer only: the counted proof cites this statement by \textbackslash{}ref as a standing input; the deduction the counted proof has to make is carried out inside the counted span, and no part of that deduction is delegated to this pointer.\par\medskip
\ifdefined\hypertarget\hypertarget{sourcepointer-5}{}\fi
\noindent\textit{Source reference.} Original lemma 07FA: sources/stacks-smoothing/pinned-original-source.tex, source lines 2447--2447. Navigation pointer only: the counted proof cites this statement by \textbackslash{}ref as a standing input; the deduction the counted proof has to make is carried out inside the counted span, and no part of that deduction is delegated to this pointer.\par\medskip
\ifdefined\hypertarget\hypertarget{sourcepointer-6}{}\fi
\noindent\textit{Source reference.} Original lemma 07FD: sources/stacks-smoothing/pinned-original-source.tex, source lines 2523--2523. Navigation pointer only: the counted proof cites this statement by \textbackslash{}ref as a standing input; the deduction the counted proof has to make is carried out inside the counted span, and no part of that deduction is delegated to this pointer.\par\medskip
\ifdefined\hypertarget\hypertarget{sourcepointer-7}{}\fi
\noindent\textit{Source reference.} Original lemma 07CE: sources/stacks-smoothing/pinned-original-source.tex, source lines 559--559. Navigation pointer only: the counted proof cites this statement by \textbackslash{}ref as a standing input; the deduction the counted proof has to make is carried out inside the counted span, and no part of that deduction is delegated to this pointer.\par\medskip
\ifdefined\hypertarget\hypertarget{sourcepointer-8}{}\fi
\noindent\textit{Source reference.} Original lemma 07F8: sources/stacks-smoothing/pinned-original-source.tex, source lines 2314--2314. Navigation pointer only: the counted proof cites this statement by \textbackslash{}ref as a standing input; the deduction the counted proof has to make is carried out inside the counted span, and no part of that deduction is delegated to this pointer.\par\medskip
\ifdefined\hypertarget\hypertarget{sourcepointer-9}{}\fi
\noindent\textit{Source reference.} Original proposition 07CM: sources/stacks-smoothing/pinned-original-source.tex, source lines 1591--1591. Navigation pointer only: the counted proof cites this statement by \textbackslash{}ref as a standing input; the deduction the counted proof has to make is carried out inside the counted span, and no part of that deduction is delegated to this pointer.\par\medskip
\ifdefined\hypertarget\hypertarget{sourcepointer-10}{}\fi
\noindent\textit{Source reference.} Original reference 07F7: sources/stacks-smoothing/pinned-original-source.tex, source lines 2296--2296. Navigation pointer only: the counted proof cites this statement by \textbackslash{}ref as a standing input; the deduction the counted proof has to make is carried out inside the counted span, and no part of that deduction is delegated to this pointer.\par\medskip
\begin{lemma}
\label{lemma-resolve-special}
Let $k \to A \to \Lambda \supset \mathfrak q$ be as in
Situation \ref{situation-local} where
\begin{enumerate}
\item $k$ is a field,
\item $\Lambda$ is Noetherian,
\item $\mathfrak q$ is minimal over $\mathfrak h_A$,
\item $\Lambda_\mathfrak q$ is a regular local ring, and
\item the field extension $\kappa(\mathfrak q)/k$ is separable.
\end{enumerate}
Then $k \to A \to \Lambda \supset \mathfrak q$ can be resolved.
\end{lemma}

\begin{proof}
Set $d = \dim \Lambda_\mathfrak q$. Set $R = k[x_1, \ldots, x_d]$.
Choose $n > 0$ such that
$\mathfrak q^n\Lambda_\mathfrak q \subset \mathfrak h_A\Lambda_\mathfrak q$
which is possible as $\mathfrak q$ is minimal over $\mathfrak h_A$.
Choose generators $a_1, \ldots, a_r$ of $H_{A/R}$. Set
$$
B = A[x_1, \ldots, x_d, z_{ij}]/(x_i^n - \sum z_{ij}a_j)
$$
Each $B_{a_j}$ is smooth over $R$ it is a polynomial
algebra over $A_{a_j}[x_1, \ldots, x_d]$ and $A_{a_j}$ is smooth over $k$.
Hence $B_{x_i}$ is smooth over $R$. Let $B \to C$ be the $R$-algebra
map constructed in Lemma \ref{lemma-improve-presentation}
which comes with a $R$-algebra retraction $C \to B$. In particular
a map $C \to \Lambda$ fitting into the diagram above.
By construction $C_{x_i}$ is a smooth $R$-algebra with
$\Omega_{C_{x_i}/R}$ free. Hence we can find $c > 0$
such that $x_i^c$ is strictly standard in $C/R$, see
Lemma \ref{lemma-compare-standard}.
Now choose $\pi_1, \ldots, \pi_d \in \Lambda$ as in
Lemma \ref{lemma-find-sequence}
where $n = n$, $e = 8c$, $\mathfrak q = \mathfrak q$ and $I = \mathfrak h_A$.
Write $\pi_i^n = \sum \lambda_{ij} a_j$ for some $\pi_{ij} \in \Lambda$.
There is a map $B \to \Lambda$ given by $x_i \mapsto \pi_i$
and $z_{ij} \mapsto \lambda_{ij}$. Set $R = k[x_1, \ldots, x_d]$.
Diagram
$$
\xymatrix{
R \ar[r] & B \ar[rd] \\
k \ar[u] \ar[r] & A \ar[u] \ar[r] & \Lambda
}
$$
Now we apply
Lemma \ref{lemma-lift-solution}
to $R \to C \to \Lambda \supset \mathfrak q$
and the sequence of elements $x_1^c, \ldots, x_d^c$ of $R$.
Assumption (2) is clear. Assumption (1) holds for $R$
by inspection and for $\Lambda$ by our choice of
$\pi_1, \ldots, \pi_d$. (Note that if
$\text{Ann}_\Lambda(\pi) = \text{Ann}_\Lambda(\pi^2)$, then we have
$\text{Ann}_\Lambda(\pi) = \text{Ann}_\Lambda(\pi^c)$ for all $c > 0$.)
Thus it suffices to resolve
$$
R/(x_1^e, \ldots, x_d^e) \to C/(x_1^e, \ldots, x_d^e) \to
\Lambda/(\pi_1^e, \ldots, \pi_d^e) \supset
\mathfrak q/(\pi_1^e, \ldots, \pi_d^e)
$$
for $e = 8c$. By
Lemma \ref{lemma-delocalize-height-zero}
it suffices to resolve this after localizing at $\mathfrak q$.
But since $x_1, \ldots, x_d$ map to a regular sequence
in $\Lambda_\mathfrak q$ we see that $R_\mathfrak p \to \Lambda_\mathfrak q$
is flat, see Algebra, Lemma \ref{algebra-lemma-flat-over-regular}. Hence
$$
R_\mathfrak p/(x_1^e, \ldots, x_d^e) \to
\Lambda_\mathfrak q/(\pi_1^e, \ldots, \pi_d^e)
$$
is a flat ring map of Artinian local rings.
Moreover, this map induces a separable field extension
on residue fields by assumption. Thus this map is a filtered colimit
of smooth algebras by
Algebra, Lemma \ref{algebra-lemma-colimit-syntomic}
and Proposition \ref{proposition-lift}.
Existence of the desired solution follows from
Algebra, Lemma \ref{algebra-lemma-when-colimit}.
\end{proof}

\end{document}
